%-------------------------
% Resume in Latex
% Author : Jake Gutierrez
% Based off of: https://github.com/sb2nov/resume
% License : MIT
%------------------------

\documentclass[letterpaper,11pt]{article}

\usepackage{amsmath}
\usepackage{latexsym}
\usepackage[empty]{fullpage}
\usepackage{titlesec}
\usepackage{marvosym}
\usepackage[usenames,dvipsnames]{color}
\usepackage{verbatim}
\usepackage{enumitem}
\usepackage[hidelinks]{hyperref}
\usepackage{fancyhdr}
\usepackage[english]{babel}
\usepackage{tabularx}
\input{glyphtounicode}

\pagestyle{fancy}
\fancyhf{} 
\fancyfoot{}
\renewcommand{\headrulewidth}{0pt}
\renewcommand{\footrulewidth}{0pt}

% Adjust margins
\addtolength{\oddsidemargin}{-0.5in}
\addtolength{\evensidemargin}{-0.5in}
\addtolength{\textwidth}{1in}
\addtolength{\topmargin}{-.5in}
\addtolength{\textheight}{1.0in}

\urlstyle{same}

\raggedbottom
\raggedright
\setlength{\tabcolsep}{0in}

% Sections formatting
\titleformat{\section}{
  \vspace{-4pt}\scshape\raggedright\large
}{}{0em}{}[\color{black}\titlerule \vspace{-5pt}]

\pdfgentounicode=1

%-------------------------
% Custom commands
\newcommand{\resumeItem}[1]{
  \item\small{
    {#1 \vspace{-2pt}}
  }
}

\newcommand{\resumeSubheading}[4]{
  \vspace{-2pt}\item
    \begin{tabular*}{0.97\textwidth}[t]{l@{\extracolsep{\fill}}r}
      \textbf{#1} & #2 \\
      \textit{\small#3} & \textit{\small #4} \\
    \end{tabular*}\vspace{-7pt}
}

\newcommand{\resumeProjectHeading}[2]{
    \item
    \begin{tabular*}{0.97\textwidth}{l@{\extracolsep{\fill}}r}
      \small#1 & #2 \\
    \end{tabular*}\vspace{-7pt}
}

\newcommand{\resumeSubHeadingListStart}{\begin{itemize}[leftmargin=0.15in, label={}]}
\newcommand{\resumeSubHeadingListEnd}{\end{itemize}}
\newcommand{\resumeItemListStart}{\begin{itemize}}
\newcommand{\resumeItemListEnd}{\end{itemize}\vspace{-5pt}}

% NEW COMMAND: Use this specifically when another position follows immediately
\newcommand{\resumeItemListEndTight}{\end{itemize}\vspace{-12pt}}

\newcommand{\resumeMyPosition}[2]{
    \vspace{-9pt}\item
    \begin{tabular*}{0.97\textwidth}[t]{l@{\extracolsep{\fill}}r}
      \textit{\small#1} & \textit{\small #2} \\
    \end{tabular*}\vspace{-6pt}
}

%-------------------------------------------
%%%%%%  RESUME STARTS HERE  %%%%%%%%%%%%%%%%%%%%%%%%%%%%

\begin{document}

%----------HEADING----------
\begin{center}
    \textbf{\Huge \scshape Manan Tuteja} \\ \vspace{1pt}
    \small 424-392-0891 $|$ \href{mailto:manan.tuteja5@gmail.com}{\underline{manan.tuteja5@gmail.com}} $|$ 
    \href{https://www.linkedin.com/in/manan-tuteja-8042b0396/}{\underline{linkedin.com/in/MananTuteja}} $|$
    \href{https://mtuteja1.github.io/portfolio/}{\underline{Portfolio}}
\end{center}

% About Me
\section{Summary}
    \textnormal
    {Graduate Mechanical Engineering student specializing in robotic design and control systems. Experience bridging mechanical and electrical CAD, with basic practical programming skills.}

%-----------EXPERIENCE-----------
\section{Experience}
  \resumeSubHeadingListStart
    % \resumeSubheading{Drone (Personal Project)}{Dec. 2025 -- Present}{}{}
    %   \vspace{-17pt}
    % \resumeItemListStart
    %   \resumeItem{Developing custom autonomous drone using Arduino Uno. The drone will be capable of a \textbf{2 kg} lift capacity with a \textbf{10 minute} flight time.}
    %   \resumeItem{Designed custom \textbf{2 layer PCB} as a power distribution board (PDB) with \textbf{EasyEDA}. The PDB allows for a 3s LiPo battery to distribute power to all electronics with fuses for safety. Due to the high amperage, the 2 oz copper PCB has two layers that are stiched with vias and use copper areas.}
    % \resumeItemListEnd

    \resumeSubheading
      {Novartis Pharmaceuticals R\&D}{Jun. 2026 -- Dec. 2026}
      {Robotics Mechanical Engineering Intern}{}
      \resumeItemListStart
        \resumeItem{Development of a fully autonomous \textbf{6-DOF robot} for \textbf{pick and place} samples using a Beckhoff PLC with \textbf{TwinCAT} on \textbf{EtherCAT} protocol. Integrated custom \textbf{HMI} interface for use in a laboratory environment.}
        \resumeItem{Implemented state machine on \textbf{PLC} through state transition table and state machine diagram.}
        \resumeItem{Custom end effector design through \textbf{Creo} using \textbf{resin 3-D printer} with custom mounting for components made with \textbf{Onyx Markforged} filament. Created engineering drawings for custom milled parts using \textbf{GD\&T} and \textbf{ASME} standards.}
        \resumeItem{Developed pneumatic and electrical circuitry integrating vacuum systems, solenoid valves, regulators, gauges, peristaltic pumps with power supplies, circuit breakers, switches, and relays through Visio.}
      \resumeItemListEnd
    \resumeSubheading
      {Triton Robotics @ University of California San Diego}{Jan. 2024 -- Present}
      {Mechanical Team Lead}{Jul. 2025 -- Present}
      \resumeItemListStart
        \resumeItem{Robot developed for ARC Robotics Competition with computer vision based turret system.}
        \resumeItem{Managed a team of over 10 mechanical engineers and general management of 30 member mechanical team in \textbf{Onshape} enterprise.}
        \resumeItem{Directed development of shooting mechanism and pitch mechanism made with flywheels, spur gears, and sensors such as encoders and IMUs.}
        \resumeItem{Designed and manufactured indexer, ball path, customizable yaw mechanism, and a rectangular x-drive mecanum chassis with tangential suspension.}
        \resumeItem{Developed custom \textbf{4-layer PCB} breakout board for STM32 Nucleo-F446RE with \textbf{EasyEDA STD} including SPI, CAN Transceiver, UART, I2C, PWM, and Analog connections.}
      \resumeItemListEnd
    \vspace{5 pt}
    \resumeMyPosition{Mechanical Team Member}{Jan. 2024 -- Jun. 2025}
        \resumeItemListStart
        \resumeItem{Robot developed for Robomasters North America Competition with computer vision based turret system.}
        \resumeItem{Designed \textbf{linear suspension system} emphasizing cost-effectiveness and assembly efficiency with \textbf{Solidworks}.}
        \resumeItem{Designed \textbf{octagonal omniwheel x-drive chassis} with custom polycarbonate electronics plate.}
      \resumeItemListEnd

    \resumeSubheading
      {University of California San Diego CENG}{Sep. 2025 -- Dec. 2025}
      {CENG 15R Reader}{}
      \resumeItemListStart
        \resumeItem{Reader/Grader for Engineering Computation Using MATLAB for 150+ Chemical and Nano Engineering Students.}
      \resumeItemListEndTight
    \vspace{-5 pt} 
    % \resumeSubheading
    %   {University of California San Diego MAE}{Sep. 2024 -- Dec. 2024}
    %   {MAE 8 Reader}{}
      % \resumeItemListStart
      %   \resumeItem{Reader/Grader for MATLAB Programming for Engineering Analysis class for Mechanical and Aerospace Engineering students.}
      % \resumeItemListEnd % <--- Standard spacing preserved here
      % \resumeSubheading
      % {FIRST Robotics Team FRC 687}{Aug. 2020 -- Jun. 2023}
      % {Design Engineer}{}
      % \resumeItemListStart
      % \resumeItem{Designed \textbf{tank drive base} with VEX Falcon 500 motors with chain and sprocket.}
      % \resumeItem{Led turret manufacturing of \textbf{variable hood for pitch} control and 3D printed \textbf{PCABS spur gears}.}  
      % \resumeItem{Design, deployment, and manufacturing of 3D printed differential bevel geared wrist using Fusion 360}
      % \resumeItemListEnd % <--- Standard spacing preserved here
  \resumeSubHeadingListEnd

%-----------EDUCATION-----------
\section{Education}
  \resumeSubHeadingListStart
    \resumeSubheading
      {University of California San Diego}{La Jolla, CA}
      {M.S. Mechanical Engineering}{Sep. 2026 -- Dec. 2027}
    \resumeSubheading
      {University of California San Diego}{La Jolla, CA}
      {B.S. Aerospace Engineering}{Sep. 2023 -- Jun. 2026}
      \resumeItemListStart
      \resumeItem{Cumulative GPA: 3.818}
      \resumeItem{Major GPA: 3.750}
      \resumeItemListEnd
  \resumeSubHeadingListEnd
%SKILLS
\section{Technical Skills}
 \begin{itemize}[leftmargin=0.15in, label={}]
    \small{\item{
     \textbf{CAD}{: Creo, Onshape, Autodesk Inventor, Autodesk Fusion 360, Solidworks, GD\&T, ASME Y14.5}  \\
     \textbf{Programming}{: Matlab/Simulink, Python (pandas, NumPy, Matplotlib, OpenCV and ROS2), C++ (Arduino), RaspberryPI, PLC (TwinCAT)} \\
     \textbf{Manufacturing}{: Lasercutting, Waterjetting, Manual Mill, General Machine Shop tools} \\
     \textbf{PCB Design}{: EasyEDA} \\
    }}
 \end{itemize}

\end{document}